\chapter{Media di funzioni e varianza nel continuo}

Nel Capitolo~7 avevamo misurato la dispersione delle variabili discrete con il valore quadratico medio e la varianza,
e per calcolarli ci era servito il teorema fondamentale per la media di una funzione, $E[g(X)] = \sum_x g(x)\,
p_X(x)$. In questo capitolo portiamo tutto nel continuo: prima il teorema fondamentale, con le somme che diventano
integrali; poi il valore quadratico medio e la varianza delle variabili continue notevoli. Alla fine raccogliamo in un
«dizionario» le corrispondenze fra il mondo discreto e quello continuo.

% ============================================================
\section{La media di una funzione di variabile continua}
% ============================================================

Sia $X \sim f_X(x)$ e sia $Y = g(X)$, con $\mathcal{Y} = g(\mathcal{X})$. Per calcolare $E[Y]$ si potrebbe ricavare la
densità $f_Y$ con i metodi del capitolo precedente e poi integrare $y\, f_Y(y)$. Ma, come nel discreto, non serve.

\teorema[valoremedioc]{Teorema fondamentale per il calcolo della media}{
  Qualunque sia la funzione $g$ (purché l'integrale converga):
  \[ E[Y] = E[g(X)] = \int_{\R} g(x)\, f_X(x) \dd x . \]
}

\begin{dimostrazione}[con la quantizzazione]
\passo{Si quantizza $X$} Come per giustificare la media (Capitolo~10), si definisce una versione quantizzata di $X$:
\[ X^{\Delta} = x_i \in [i\Delta, (i+1)\Delta[ \qquad \text{se} \quad i\Delta \leq X < (i+1)\Delta . \]
\passo{Si applica $g$} La variabile $Y^{\Delta} = g(X^{\Delta})$ è discreta: vale $g(x_i)$ quando $X$ cade
nell'$i$-esimo intervallo. Per il teorema fondamentale del caso discreto (Teorema~\ref{teo:valoremedio}):
\[ E[Y^{\Delta}] = \sum_i g(x_i)\, P(X^{\Delta} = x_i) = \sum_i g(x_i)\, P(x_i \leq X < x_{i+1}) . \]
\passo{Si stringono gli intervalli} Per $\Delta$ piccolo $P(x_i \leq X < x_{i+1}) \simeq f_X(x_i)\,\Delta$, e
$Y^{\Delta}$ si avvicina a $Y$. Se $g(x) f_X(x)$ è integrabile secondo Riemann, la somma diventa un integrale:
\[ E[Y] = \lim_{\Delta \to 0} E[Y^{\Delta}] = \lim_{\Delta \to 0} \sum_i g(x_i)\, \underbrace{f_X(x_i)\,\Delta}_{\simeq
   P(x_i \leq X < x_{i+1})} = \int_{\R} g(x)\, f_X(x) \dd x . \]
\end{dimostrazione}

\begin{esempio}{Due strade per la stessa media}
Sia $X \sim \mathcal{U}(0, 1)$ e $Y = X^2$.
\begin{itemize}
  \item \textbf{Con il teorema fondamentale}: $E[Y] = \int_0^1 x^2 \cdot 1 \dd x = \left[ \frac{x^3}{3} \right]_0^1 =
        \frac13$.
  \item \textbf{Passando per la densità di $Y$}: su $(0, 1)$ la funzione $x^2$ è crescente, con $x(y) = \sqrt{y}$ e
        $\dot g(x(y)) = 2\sqrt{y}$; quindi $f_Y(y) = \frac{1}{2\sqrt{y}}$ per $0 < y < 1$, e
        \[ E[Y] = \int_0^1 y\, \frac{1}{2\sqrt{y}} \dd y = \int_0^1 \frac{\sqrt{y}}{2} \dd y = \frac12 \left[ \frac{y^{3/2}}{3/2}
           \right]_0^1 = \frac12 \cdot \frac23 = \frac13 . \]
\end{itemize}
Stesso risultato, ma la prima strada è molto più breve. (Si è usata la regola delle potenze anche con l'esponente non
intero $\frac12$: la formula $(x^n)' = n\,x^{n-1}$ vale per qualsiasi esponente reale.)
\end{esempio}

Dal teorema e dalla linearità dell'integrale seguono subito le proprietà della media già viste nel discreto:
\[ E[aX + b] = \int (ax + b)\, f_X(x) \dd x = a \int x\, f_X(x) \dd x + b \int f_X(x) \dd x = a\, E[X] + b , \]
e più in generale $E[g_1(X) + g_2(X)] = E[g_1(X)] + E[g_2(X)]$.

% ============================================================
\section{Valore quadratico medio e varianza}
% ============================================================

A questo punto la generalizzazione dei concetti introdotti per le variabili discrete è immediata. Data una variabile
aleatoria $X \sim f_X(x)$, con media $\mu_X = E[X]$, si definiscono:
\begin{itemize}
  \item il \textbf{valore quadratico medio} (\emph{mean square}): $\displaystyle X_{rms}^2 = E[X^2] = \int_{\R} x^2\,
        f_X(x) \dd x$;
  \item il \textbf{valore efficace} (\emph{root mean square}): $\displaystyle X_{rms} = \sqrt{E[X^2]} = \sqrt{\int_{\R}
        x^2\, f_X(x) \dd x}$;
  \item la \textbf{varianza}: $\displaystyle \sigma_X^2 = E\left[(X - \mu_X)^2\right] = \int_{\R} (x - \mu_X)^2\, f_X(x)
        \dd x$;
  \item la \textbf{deviazione standard}: $\sigma_X = \sqrt{\sigma_X^2}$.
\end{itemize}

\teorema[varianzac]{Varianza e valore quadratico medio}{
  \[ \sigma_X^2 = E[X^2] - \mu_X^2 = X_{rms}^2 - \mu_X^2 , \qquad \sigma_X = \sqrt{E[X^2] - \mu_X^2} . \]
}

\begin{dimostrazione}
\passo{} Si sviluppa il quadrato e si usa la linearità della media (le costanti $\mu_X$ e $\mu_X^2$ escono):
\[ \sigma_X^2 = E\left[ X^2 + \mu_X^2 - 2\mu_X X \right] = E[X^2] + \mu_X^2 - 2\mu_X\, E[X] = E[X^2] + \mu_X^2 - 2\mu_X^2
   = E[X^2] - \mu_X^2 . \]
\end{dimostrazione}

Tutte le proprietà della media e della varianza viste per le variabili aleatorie discrete valgono, ovviamente, anche per
quelle continue: le dimostrazioni sono le stesse, con gli integrali al posto delle somme. In particolare
$\sigma_{aX + b}^2 = a^2 \sigma_X^2$, e valgono le disuguaglianze di Markov e di Chebyshev. Per esempio, per una variabile
$Z \geq 0$ con densità $f_Z$ la dimostrazione di Markov diventa
\[ P(Z \geq \delta) = \int_{\delta}^{\infty} f_Z(z) \dd z \leq \int_{\delta}^{\infty} \frac{z}{\delta}\, f_Z(z) \dd z \leq
   \int_0^{\infty} \frac{z}{\delta}\, f_Z(z) \dd z = \frac{E[Z]}{\delta} , \]
con gli stessi tre passaggi del Capitolo~7: si moltiplica per $\frac{z}{\delta} \geq 1$, si estende l'integrale aggiungendo
una parte non negativa, si riconosce la media.

% ============================================================
\section{Le varianze delle variabili notevoli}
% ============================================================

\subsection{La variabile uniforme}

\teorema[varuniformec]{Varianza dell'uniforme}{
  Se $X \sim \mathcal{U}(a, b)$:
  \[ E[X] = \frac{a + b}{2} , \qquad E[X^2] = \frac{a^2 + ab + b^2}{3} , \qquad \sigma_X^2 = \frac{(b - a)^2}{12} ,
     \qquad \sigma_X = \frac{b - a}{\sqrt{12}} . \]
}

\begin{dimostrazione}
\passo{La media} È già stata calcolata nel Capitolo~11.
\passo{Il valore quadratico medio} Con la primitiva $\frac{x^3}{3}$:
\[ E[X^2] = \frac{1}{b - a} \int_a^b x^2 \dd x = \frac{1}{b - a} \left[ \frac{x^3}{3} \right]_a^b = \frac{b^3 - a^3}{3(b -
   a)} = \frac{(b - a)(b^2 + ab + a^2)}{3(b - a)} = \frac{b^2 + ab + a^2}{3} , \]
dove si è usato il prodotto notevole $b^3 - a^3 = (b - a)(b^2 + ab + a^2)$ (basta moltiplicare a destra per verificarlo:
$b^3 + ab^2 + a^2 b - ab^2 - a^2 b - a^3 = b^3 - a^3$).
\passo{La varianza} Al denominatore comune 12:
\[
  \sigma_X^2 = \frac{b^2 + ab + a^2}{3} - \frac{(b + a)^2}{4} = \frac{4b^2 + 4ab + 4a^2 - 3b^2 - 3a^2 - 6ab}{12}
  = \frac{b^2 - 2ab + a^2}{12} = \frac{(b - a)^2}{12} .
\]
\end{dimostrazione}

La varianza dipende solo dalla \textbf{larghezza} dell'intervallo, non dalla sua posizione: spostare l'intervallo sposta
la media ma non la dispersione. Si confronti con l'uniforme discreta su $M$ valori (Capitolo~7), che ha varianza
$\frac{M^2 - 1}{12}$: per $M$ grande è circa $\frac{M^2}{12}$, cioè (larghezza)$^2/12$, lo stesso risultato.

\begin{esempio}{L'errore di quantizzazione}
In un convertitore A/D con passo $\Delta$ (Capitolo~12), l'errore che si commette sostituendo $X$ con il centro del
suo intervallo sta fra $-\frac{\Delta}{2}$ e $\frac{\Delta}{2}$; se gli intervalli sono stretti rispetto alle variazioni
della densità, l'errore è circa uniforme su quell'intervallo. La sua potenza (il valore quadratico medio, che coincide
con la varianza perché la media è 0) è allora
\[ \sigma^2 = \frac{\Delta^2}{12} . \]
È la celebre formula del \textbf{rumore di quantizzazione}: ogni bit in più dimezza $\Delta$ e divide la potenza
dell'errore per 4.
\end{esempio}

\subsection{La variabile esponenziale}

\teorema[varesponenziale]{Varianza dell'esponenziale}{
  Se $X \sim \mathcal{E}(\lambda)$:
  \[ E[X] = \frac{1}{\lambda} , \qquad E[X^2] = \frac{2}{\lambda^2} , \qquad \sigma_X^2 = \frac{1}{\lambda^2} , \qquad
     \sigma_X = \frac{1}{\lambda} . \]
}

\begin{dimostrazione}
\passo{Il valore quadratico medio, per parti} Il gradino riduce l'integrale a $[0, \infty)$. Con $f(x) = x^2$ (derivata
$2x$) e $g'(x) = \lambda e^{-\lambda x}$ (primitiva $-e^{-\lambda x}$):
\[ E[X^2] = \int_0^{\infty} x^2 \lambda e^{-\lambda x} \dd x = \Big[ x^2 \left( -e^{-\lambda x} \right) \Big]_0^{\infty} -
   \int_0^{\infty} 2x \left( -e^{-\lambda x} \right) \dd x . \]
\passo{Il bordo è nullo} In 0 vale 0, e $x^2 e^{-\lambda x} \to 0$ per $x \to \infty$ (l'esponenziale batte la potenza).
\passo{Si riconosce la media} Moltiplicando e dividendo per $\lambda$:
\[ E[X^2] = 2 \int_0^{\infty} x\, e^{-\lambda x} \dd x = \frac{2}{\lambda} \underbrace{\int_0^{\infty} x\, \lambda e^{-\lambda
   x} \dd x}_{E[X] = \frac{1}{\lambda}} = \frac{2}{\lambda^2} . \]
\passo{La varianza} $\sigma_X^2 = \frac{2}{\lambda^2} - \frac{1}{\lambda^2} = \frac{1}{\lambda^2}$.
\end{dimostrazione}

Per l'esponenziale la deviazione standard è uguale alla media: un tempo di attesa medio di 10 minuti fluttua, tipicamente,
di altri 10 minuti. È una variabile molto «aleatoria», con rapporto $\mu_X/\sigma_X = 1$.

\subsection{La variabile di Laplace}

\teorema[varlaplace]{Varianza della variabile di Laplace}{
  Se $X \sim \mathcal{L}(\lambda)$:
  \[ E[X] = 0 , \qquad E[X^2] = \sigma_X^2 = \frac{2}{\lambda^2} , \qquad \sigma_X = \frac{\sqrt{2}}{\lambda} . \]
}

\begin{dimostrazione}
\passo{La parità} $x^2 \frac{\lambda}{2} e^{-\lambda\abs{x}}$ è pari: l'area a sinistra dello zero è uguale a quella a
destra, e l'integrale su $\R$ è il doppio di quello su $[0, \infty)$:
\[ E[X^2] = \int_{-\infty}^{\infty} x^2\, \frac{\lambda}{2}\, e^{-\lambda\abs{x}} \dd x = 2 \cdot \frac{\lambda}{2}
   \int_0^{\infty} x^2 e^{-\lambda x} \dd x = \int_0^{\infty} x^2\, \lambda e^{-\lambda x} \dd x = \frac{2}{\lambda^2} . \]
L'ultimo integrale è il valore quadratico medio dell'esponenziale unilatera, calcolato sopra.
\passo{La varianza} La media è nulla, quindi $\sigma_X^2 = E[X^2] - 0 = \frac{2}{\lambda^2}$.
\end{dimostrazione}

\subsection{La variabile di Cauchy}

Per la variabile di Cauchy $\mathcal{C}(a, b)$ non esistono né la media né la varianza, e quindi nemmeno il valore
efficace e la deviazione standard. Per la media lo abbiamo visto nel Capitolo~11. Il valore quadratico medio è infinito:
con $a = 0$ e $b = 1$, l'integrando $\frac{x^2}{\pi(1 + x^2)}$ tende a $\frac{1}{\pi}$ per $x \to \pm\infty$, e una funzione
che non scende sotto un valore positivo ha area infinita su una semiretta (come il rettangolo infinito del
Capitolo~10). Senza media, poi, la varianza non si può nemmeno definire.

\info{Che cosa ci dice la deviazione standard, nel continuo}{
  La disuguaglianza di Chebyshev garantisce che la probabilità di allontanarsi dalla media di almeno $k$ deviazioni
  standard sia al più $\frac{1}{k^2}$, per qualsiasi densità. La tabella confronta il limite con i valori veri per tre
  variabili notevoli (per la gaussiana, la colonna verrà spiegata nel Capitolo~15).
  \begin{center}
  \renewcommand{\arraystretch}{1.25}
  \begin{tabular}{@{}c|cccc|c@{}}
    $k$ & uniforme & esponenziale & Laplace & gaussiana & Chebyshev $\frac{1}{k^2}$ \\ \hline
    1 & 0,423 & 0,135 & 0,243 & 0,317 & 1 \\
    1,5 & 0,134 & 0,082 & 0,120 & 0,134 & 0,444 \\
    2 & 0 & 0,050 & 0,059 & 0,046 & 0,25 \\
    3 & 0 & 0,018 & 0,014 & 0,003 & 0,111
  \end{tabular}
  \end{center}
  Per esempio, per l'esponenziale con $\lambda = 1$ ($\mu = \sigma = 1$) e $k = 2$:
  $P(\abs{X - 1} \geq 2) = P(X \geq 3) = e^{-3} \approx 0{,}050$. L'uniforme non esce mai oltre $1{,}73\,\sigma$ dal centro
  (la semilarghezza dell'intervallo è $\frac{b - a}{2} = \sqrt{3}\,\sigma$).
}

% ============================================================
\section{Il dizionario fra discreto e continuo}
% ============================================================

Ora che abbiamo visto entrambi i mondi, conviene metterli uno accanto all'altro: quasi tutte le formule si traducono
con poche regole.

\begin{table}[H]
\centering
\renewcommand{\arraystretch}{1.6}
\begin{tabular}{@{}lll@{}}
  \toprule
  \textbf{Concetto} & \textbf{Variabile discreta} & \textbf{Variabile continua} \\
  \midrule
  descrizione completa & pmf $p_X(x) = P(X = x)$ & densità $f_X(x)$, con $P(X = x) = 0$ \\
  probabilità di un insieme & $\sum_{x \in A} p_X(x)$ & $\int_A f_X(x) \dd x$ \\
  normalizzazione & $\sum_x p_X(x) = 1$ & $\int_{\R} f_X(x) \dd x = 1$ \\
  CDF $F_X(x) = P(X \leq x)$ & a scala, salti alti $p_X(x)$ & continua, $F_X' = f_X$ \\
  media & $\sum_x x\, p_X(x)$ & $\int x\, f_X(x) \dd x$ \\
  media di $g(X)$ & $\sum_x g(x)\, p_X(x)$ & $\int g(x)\, f_X(x) \dd x$ \\
  funzione $Y = g(X)$ & $p_Y(y) = \sum_{g(x) = y} p_X(x)$ & $f_Y(y) = \sum_{g(x_i) = y} \frac{f_X(x_i)}{\abs{\dot g(x_i)}}$ \\
  condizionamento a $A$ & $p_{X \mid A}(x) = \frac{P(\{X = x\} \cap A)}{P(A)}$ & $F_{X \mid A}(x) = \frac{P(\{X \leq x\}
    \cap A)}{P(A)}$ \\
  varianza & $E[X^2] - \mu_X^2$ & $E[X^2] - \mu_X^2$ \\
  \bottomrule
\end{tabular}
\caption{Le regole di traduzione: le somme diventano integrali, $p_X(x)$ diventa $f_X(x) \dd x$; l'unica formula davvero
nuova è quella delle funzioni di variabile aleatoria, con la divisione per $\abs{\dot g}$.}
\end{table}

\begin{table}[H]
\centering
\renewcommand{\arraystretch}{1.6}
\begin{tabular}{@{}lllll@{}}
  \toprule
  \textbf{Variabile} & \textbf{Densità} & \textbf{Media} & \textbf{Valore quadratico medio} & \textbf{Varianza} \\
  \midrule
  uniforme $\mathcal{U}(a, b)$ & $\frac{1}{b - a}$ su $(a, b)$ & $\frac{a + b}{2}$ & $\frac{a^2 + ab + b^2}{3}$ &
    $\frac{(b - a)^2}{12}$ \\
  esponenziale $\mathcal{E}(\lambda)$ & $\lambda e^{-\lambda x} u(x)$ & $\frac{1}{\lambda}$ & $\frac{2}{\lambda^2}$ &
    $\frac{1}{\lambda^2}$ \\
  Laplace $\mathcal{L}(\lambda)$ & $\frac{\lambda}{2} e^{-\lambda\abs{x}}$ & 0 & $\frac{2}{\lambda^2}$ & $\frac{2}{\lambda^2}$ \\
  Cauchy $\mathcal{C}(a, b)$ & $\frac{1}{b\pi} \frac{1}{1 + ((x - a)/b)^2}$ & non esiste & $+\infty$ & non esiste \\
  \bottomrule
\end{tabular}
\caption{Medie e varianze delle variabili continue notevoli.}
\end{table}

\gbox{In sintesi}{mspiA!80!black}{
\begin{itemize}
  \item \textbf{Teorema fondamentale}: $E[g(X)] = \int g(x)\, f_X(x) \dd x$, senza bisogno della densità di $g(X)$; si
        dimostra quantizzando $X$ e usando il caso discreto.
  \item Valore quadratico medio $\int x^2 f_X$, valore efficace, varianza $\int (x - \mu_X)^2 f_X = E[X^2] - \mu_X^2$,
        deviazione standard: le definizioni e le proprietà sono quelle del discreto, con gli integrali.
  \item Uniforme: $\sigma^2 = \frac{(b-a)^2}{12}$ (errore di quantizzazione $\frac{\Delta^2}{12}$). Esponenziale:
        $\sigma^2 = \frac{1}{\lambda^2}$, $\sigma = \mu$. Laplace: $\sigma^2 = \frac{2}{\lambda^2}$. Cauchy: né media né
        varianza.
  \item Markov e Chebyshev valgono anche nel continuo; per le densità notevoli le probabilità vere sono molto più
        piccole del limite $\frac{1}{k^2}$.
\end{itemize}
}
